\documentclass[letterpaper,11pt]{article}
\usepackage[margin=0.66in,top=0.72in,bottom=0.68in,headheight=15pt,headsep=15pt]{geometry}
\usepackage[T1]{fontenc}
\usepackage[scaled=.94]{helvet}
\renewcommand{\familydefault}{\sfdefault}
\usepackage{tikz,amsmath,amssymb,array,tabularx,fancyhdr,microtype,xcolor,enumitem}
\usetikzlibrary{arrows.meta}
\usepackage[hidelinks]{hyperref}
\definecolor{blue}{RGB}{20,73,112}
\pagestyle{fancy}\fancyhf{}\renewcommand{\headrulewidth}{0pt}\renewcommand{\footrulewidth}{0pt}
\fancyhead[L]{\small Week 52 / Hinged frames and braces / Adult guide}
\fancyfoot[L]{\scriptsize Bellingham Math Circle / Week 52 / F52-FAC-v2}\fancyfoot[R]{\scriptsize\thepage}
\setlength{\emergencystretch}{2em}
\setlength{\parindent}{0pt}\setlength{\parskip}{6pt}
\setlist[itemize]{leftmargin=16pt,itemsep=2pt,topsep=3pt}
\setlist[enumerate]{leftmargin=18pt,itemsep=2pt,topsep=3pt}
\newcommand{\titleline}[1]{\par{\large\bfseries #1}\par}
\newcommand{\sub}[1]{\par\textbf{#1}\par}
\newcommand{\task}[3]{\sub{Problem #1 / student page #2 / #3}}
\newcommand{\hints}[3]{\textbf{Hints, in order:} (1) #1 (2) #2 (3) #3\par}
\newcommand{\grid}[4]{\begin{tikzpicture}[x=#3mm,y=#3mm,baseline=(current bounding box.center)]
\foreach \x in {0,...,#2}{\draw[gray] (\x,0)--(\x,#1);}
\foreach \y in {0,...,#1}{\draw[gray] (0,\y)--(#2,\y);}
\foreach \r/\c in {#4}{\draw[blue,line width=1.6pt] ({\c-1},{#1-\r})--(\c,{#1-\r+1});}
\foreach \x in {0,...,#2}{\foreach \y in {0,...,#1}{\filldraw[fill=white] (\x,\y) circle(1.6pt);}}
\foreach \r in {1,...,#1}{\node[anchor=east,font=\scriptsize] at (-.1,{#1-\r+.5}) {R\r};}
\foreach \c in {1,...,#2}{\node[anchor=south,font=\scriptsize] at ({\c-.5},{#1+.08}) {C\c};}
\end{tikzpicture}}
\newcommand{\links}[3]{\begin{tikzpicture}[x=1cm,y=.65cm,baseline=(current bounding box.center)]
\foreach \r in {1,...,#1}{\coordinate(R\r) at (0,{#2-\r});}
\foreach \c in {1,...,#2}{\coordinate(C\c) at (2.4,{#2-\c});}
\foreach \r/\c in {#3}{\draw[blue,line width=1pt](R\r)--(C\c);}
\foreach \r in {1,...,#1}{\filldraw[fill=white](R\r)circle(2pt);\node[anchor=east,font=\scriptsize]at(-.1,{#2-\r}){R\r};}
\foreach \c in {1,...,#2}{\filldraw[fill=white](C\c)circle(2pt);\node[anchor=west,font=\scriptsize]at(2.5,{#2-\c}){C\c};}
\end{tikzpicture}}
\begin{document}
\titleline{Mathematical overview: what actually holds}
\textbf{A fixed-bar triangle is rigid in the plane. A hinged square is flexible; one rigid diagonal makes it rigid.} Bars keep their lengths, pins turn without sliding, and all joints stay in one plane. Moving the entire frame is allowed and does not change its shape. A cardboard polygon with fixed corners is a different object. A braced quadrilateral can fold out of the plane; this session makes no general three-dimensional rigidity claim.

\textbf{Complete square-grid theorem.} For an initially square, complete $m$-by-$n$ cell grid, make one dot for every row strip and every column strip. A brace in cell $(i,j)$ makes the link $R_iC_j$. The grid is \emph{infinitesimally rigid exactly when this link graph is connected}. With $k$ connected groups, there are $k-1$ nontrivial first-order shape motions, after removing whole-frame motions. This is a statement about ideal bar constraints, not about stiffness under a child's push.

\textbf{Here connectedness also decides finite continuous motion, by a separate geometric argument.} Near the square placement, each four-bar cell remains a parallelogram. All vertical bars across a row strip share a direction, and all horizontal bars across a column strip share a direction. A square-length diagonal fixes those two directions at a right angle. Linked right-angle relations propagate through a connected group. If all dots are linked, every horizontal strip has one direction and every vertical strip its perpendicular: only a whole-frame turn/slide remains. If the groups are separate, rotate each group's strip directions by a different small angle and place joints by sums of the horizontal and vertical vectors. All side and braced-diagonal lengths remain fixed while an unbraced cell changes angle. This constructs a real flex; it does not assume that every infinitesimal flex can be integrated. See page 9 for the formulas.

\textbf{The minimum is $m+n-1$ braces.} There are $m+n$ initially separate dots, and one link joins at most two groups. A spanning tree attains the bound. Having that many braces is insufficient unless the dots are connected; neither touching every strip nor a cycle elsewhere repairs a separate group. A link on a cycle can be removed without disconnecting a connected graph; a bridge cannot.

\textbf{Scope.} Every side bar is present, every cell starts nondegenerate and square, braces are rigid cell diagonals of length $L\sqrt2$, and at most one is used in a cell. Keep all row/column dots, even unlinked ones. Link crossings are not joints. Do not apply this criterion to grids with holes or missing bars, strings, sliding/locked joints, long diagonals, or lifted frames. The grid theorem concerns continuous local motion from the initial nondegenerate square placement, not all configurations reached after removing a brace and passing through collapse, distant embeddings or mechanical strength.

\sub{What children are meant to discover}
K--1, Problems 1--3: a triangle resists a shape change; a square can shear; a chosen diagonal constrains it. Grades 2--3, Problems 4--5: placement matters, and three braces suffice in a 2-by-2 grid. Grades 3--5, Problems 6--11: physical contrasts acquire a useful link picture after the handoff on page 4. Ready Grades 4--5, Problems 12--14: minimum and universal claims need an argument, not only examples. Grade labels are approximate; no algebra or calculus is required. A hand test is an \emph{experiment}; a proposed rule is a \emph{conjecture}; propagation or a length-preserving construction is an \emph{explanation}.

\newpage
\titleline{Materials and preparation for eleven children}
Three anchored adults: parent with KK11; mathematician with four third graders; organizer with 445. Two pairs at each younger table; at the upper table rotate builder, tester and recorder after each design. The spare upper kit permits two designs to stay assembled at once.

\begin{tabularx}{\linewidth}{@{}Xrrrr@{}}\textbf{Preassembled kit} &\textbf{Number}&\textbf{Sides}&\textbf{Braces}&\textbf{Pins}\\\hline
K--1: separate triangle and square &2&14&2&14\\
Grades 2--3: 2-by-2 cell grid &2&24&8&18\\
Upper: 2-by-3 cell grid, including spare &2&34&12&24\\
Common demonstration square &1&4&1&4\\\hline
\textbf{Base totals} &&\textbf{76}&\textbf{23}&\textbf{60}\\
Loose spares &&8&4&12\\
\textbf{Prepare altogether} &&\textbf{84}&\textbf{27}&\textbf{72}
\end{tabularx}

Each K--1 kit uses 7 side bars, 1 brace and 7 pins. A 2-by-2 uses 12 sides, 4 braces and 9 pins. A 2-by-3 uses 17 sides, 6 braces and 12 pins. Braces reuse existing corner pins; they do not add pins at crossings. Count 72 \emph{complete} pin fasteners, including backs/washers if the chosen system needs them. The main route needs no 3-by-3 kit. Optional direct testing of paper 3-by-3 cases adds 24 sides, 9 braces and 16 pins, separately rehearsed.

\sub{Fabrication and fit}
Use stiff straight prepunched strip bars with freely turning removable pins, or a compatible reusable hinged-bar kit. Proposed strips are 14 mm wide, with holes centered 8 mm from each end. Side hole spacing is 60 mm, total length 76 mm. Diagonal hole spacing is $60\sqrt2=84.853$ mm, total length 100.853 mm. These are \emph{hole-center} dimensions. Do not substitute a bar of total length 84.85 mm. The center spans are 60-by-60 mm for a square, 120-by-120 mm for a 2-by-2 and 180-by-120 mm for a 2-by-3; strip ends extend beyond them. Student diagrams are smaller recording maps, not full-size assembly mats.

Before assembly, cut all strips using one measured side template and one diagonal template; transfer/punch both holes from the same template, then check the actual hole-center distances. Use stiff sheet stock or rigid kit bars that resist bending under gentle testing; soft single-layer card is a poor proxy. Match hole size to pins to permit rotation without appreciable sliding. At an interior joint allow four side strips and potentially four incident braces; pins must accommodate that stack without clamping it. Adults cut/punch and service fasteners. Children choose the brace cell and the push.

\textbf{Preparation forecast:} 45--60 minutes if suitable straight strip stock/fasteners and bulk cutting/punching are already available: about 20 minutes cutting/punching, 20--25 assembling, 10--15 checking/resetting. Raw-material sourcing or drilling a rigid stock can take longer; use an earlier kit-making evening. These times are estimates, not measured. Later reset/check is about 10--15 minutes. Do not spend children's exploration time fabricating kits.

Also prepare 13 pencils (11 plus 2 spares), 6 blue/red marker pairs (5 work groups plus spare), and 15 removable label tabs: $R1,R2,C1,C2,C3$ for each upper model plus a spare set. Supply 5 uncluttered work areas roughly 30-by-25 cm, one per pair/trio; spare upper model sits alongside. Keep loose pins in adult-controlled trays.

\newpage
\titleline{Printing, physical pretest and the hour}
\sub{Print from the current 12-page student packet, F52-S-v2}
Print single-sided, US Letter, 100\% scale. K--1: 2 copies each of pages 1--2, one set per pair (4 sheets). Grades 2--3: 2 copies each of pages 3--4 (4 sheets). Upper: 1 copy each of pages 5--9 (5 sheets), shared by the trio; give one at a time. Keep 1 copy each of pages 10--12 for a ready continuation/return visit (3 sheets). Total 16 student sheets. Print 3 copies of this adult guide. The spare frame does not require an extra worksheet set; its brace map can remain on the same page.

\sub{Required physical pretest: unperformed}
Check every assembled unbraced square/grid changes shape gently while flat; each triangle does not. Fit the same diagonal into a square and check that it restricts shear without forcing the holes. Check P4A/B flex and P4C/D hold, then P7A/B on the upper kits. Try the two-group handoff and every spare bar/pin. If friction makes an unbraced frame seem rigid, loosen/service pins; if a braced frame bends, slides or has slack, replace the material or diagonal. A diagram or code check does not certify these operations. Physical fabrication, fit, pin friction/slop, adult servicing workload and classroom piloting remain \textbf{unperformed}.

\sub{A flexible 60-minute route}
\begin{tabularx}{\linewidth}{@{}lX@{}}0--5 & Running around, then settle at fixed tables.\\
5--8 & Children handle their preassembled triangle/square/grid, turning whole objects and trying gentle flat pushes. No general rule yet.\\
8--11 & Whole-group legal-attempt launch below.\\
11--30 & K--1 P1--3; third graders P4 then P5; upper P6 then physical P7. Adults read/service pins without selecting cells.\\
30--48 & Stay with fruitful concrete designs. Ready upper children do the page-4 handoff, then choose P8 or P9--10. P11 is optional. Third graders may try P6 with an available upper kit after its group finishes.\\
48--55 & Keep working or use a fallback. P12--14 normally wait for a return visit; a ready group may start one.\\
55--60 & Each table shares one changed shape or brace design. Ask: ``What did you try, and what makes you think your claim holds?'' Identify experiment, conjecture or explanation.
\end{tabularx}

\sub{Launch: materials first, then one legal attempt}
Place the demonstration square flat. Say: ``Keep the bars straight and the frame on the table. Can its shape change?'' Invite a child to hold one side lightly and choose a gentle sideways push. Turn/slide the whole frame once: ``That moved it; did its shape change?'' Return to a square. Let the child choose either pair of opposite corners; the adult opens those two pins and fits the one diagonal. The child tests it. Say: ``This diagonal is a brace. Choose where braces go, then test and mark your design.'' Stop there; do not announce the grid criterion.

\textbf{Fallbacks that retain choices.} K--1: one child makes a rhombus, partner copies it with the other square; swap. Third graders: one pair invents a 2-by-2 brace design; the other predicts then tests it. Upper: design a 2-by-3 frame that holds using four braces; partner chooses a removal and predicts its effect. Photograph/mark the reference before changing it. If children lose the two-label meaning, continue physical work; adults do not operate the investigation for them.

\newpage
\titleline{The handoff: physical bar directions to links}
Do this at the upper table after physical P7 and before using paper-only P8/11--14. P7's contrast has an isolated column; children need a covered-but-separated contrast too. Use the spare 2-by-3 kit with braces $(1,1),(1,2),(2,3)$. This avoids moving the anchored adults or borrowing another table's active kit. Alternatively borrow a finished 2-by-2 P4B model: $(1,1),(2,2)$.

\begin{center}
\grid{2}{3}{17}{1/1,1/2,2/3}\hspace{15mm}\links{2}{3}{1/1,1/2,2/3}
\end{center}
The 2-by-3 model has two groups: $\{R1,C1,C2\}$ and $\{R2,C3\}$. Every strip has a brace, but the groups can turn differently. Any empty cell, $(1,3),(2,1),(2,2)$, joins them. In P4B the groups are $\{R1,C1\}$ and $\{R2,C2\}$; either empty cell joins them.

\begin{enumerate}
\item Keep the brace record. Put R1/R2 tabs beside the \emph{vertical} bars across the row strips, and C1/C2/C3 beside the \emph{horizontal} bars across the column strips. A row dot stands for a common vertical bar direction; it is not a physical pin. A column dot stands for a common horizontal direction.
\item Children choose gentle flat pushes. Ask: ``Which bars stay parallel? Which cells stay square?'' Supply this fact if needed: equal opposite side lengths keep each nondegenerate cell on its parallelogram branch near the starting square, so its opposite bars stay parallel. Propagate that across each strip. A diagonal keeps its braced cell's two directions perpendicular.
\item Use student page 6's non-task example before asking for any conversion: the 1-by-3 input has its only brace in the middle cell; the intermediate identifies R1 and C2; the output is the R1--C2 link with C1 and C3 still present. Draw the same link for one brace in this model, then let children complete the picture. Links may cross without joining.
\item Ask children to identify the two groups and find a push that lets their bar directions turn relative to one another. Let a child choose an empty cell for one added brace and predict what it couples. Adult services the pin; child tests the prediction, keeps the new map and explains what changed.
\item Give the geometric reason: linked right-angle relations propagate within a group; separate groups can have separate small turns. Ask children to point to the physical directions represented by both ends of a chosen link. Only then use links to predict a new paper frame.
\end{enumerate}

\textbf{Readiness check.} Can children distinguish a cell address from a bar direction, retain every dot, and tell a tested observation from a prediction or a fixed-length explanation? If yes, they can choose P8 or P9--11. If they can only transcribe cell addresses, return to P4--7 and preserve their map for another day. Adult reading and pin assistance are ordinary support; adult choosing every placement and tracing every group would remove the mathematical choices. This route is proposed and untested.

\textbf{No masks or frozen portions.} All side bars remain in the complete grid. Do not hold a region fixed with tape, remove cells, cover dots, or apply the criterion to a modified frame. Holding one side lightly removes whole-frame movement during a test; it must not lock internal angles. The optional additional 3-by-3 model supplies an experiment after its own rehearsal, not a substitute for the meaning of links.

\newpage
\titleline{K--1 and Grades 2--3: concrete answers}
\task{1}{1}{K--1 core}
Children try a triangle and square and draw actual new shapes. The triangle's three sides/joints are fixed in shape; the square's four sides/joints can shear into non-square rhombi. Reflections or turning a triangle are not continuous shape changes from its fixed-length placement. A child explanation is: ``The third corner has to be the same distance from both ends of this side.'' Adults can show the two possible circle intersections; there is no continuous path between them with the lengths fixed.
\hints{Try holding one side lightly on the table and push a different corner.}{Check whether the bars bent or the whole object just moved.}{Compare the square with a narrow rhombus made from its same four bars.}

\task{2}{2}{K--1 core}
One diagonal is the minimum. Either opposite-corner choice works. Zero allows shear; one splits the square into two fixed triangles. The child selects corners and records the brace. Adults attach it without squeezing side pins tight.
\hints{Point to the corners you want to join.}{Try one diagonal and the same flat push.}{Where are the two triangles now?}

\task{3}{2}{K--1 core or continuation}
Keep \emph{that same} removed diagonal. On the ordinary nondegenerate simple rhombus branch, only a square fits it. A 60-degree rhombus with 60 mm sides has diagonals 60 and $60\sqrt3\approx103.92$ mm; neither is 84.85 mm.

\textbf{The full ideal answer also allows overlap.} Call the diagonal ends $A,C$. Each other corner $B,D$ must be one of two circle intersections, distance $L$ from both ends. Distinct choices give the square; coincident choices give a doubled right-isosceles triangle with $B=D$. The unbraced frame can reach it through a fully collapsed placement. The student question forbids neither overlap nor coincident joints: credit this noticed exception. A real kit may physically prevent it; do not force pins or substitute a different bar.
\hints{Try to fit the removed bar into your new shape without bending anything.}{Compare a very thin rhombus with a shape close to square.}{Where are the four corners when both holes fit?}
If drawing is a bottleneck, an adult draws the child's chosen shape. Stay with physical copying/refitting instead of adding graph work. Optional deeper question: can a different diagonal hold a different rhombus? It can, with its appropriate length; that is a new kit/model, outside the square-grid criterion.

\task{4}{3}{Grades 2--3 core}
Rows count from the top; columns from the left. Let each pair build/test all four contrasting maps. Circle \textbf{C and D}. A has only $(1,1)$ and flexes. B has $(1,1),(2,2)$ and flexes despite a brace in every strip. C has $(1,1),(1,2),(2,1)$ and holds. D has all four and holds; its fourth brace is redundant. Draw a changed shape only for A/B. For B, two braced squares can change their relative direction through the unbraced cells; no braced square itself has to shear.
\hints{Choose a gentle push on a corner away from a braced cell.}{Can two square cells turn relative to one another?}{Try C after B using the same model and push; keep the old record.}

\task{5}{4}{Grades 2--3 core; several designs can occupy the hour}
Minimum \textbf{three}. All four minimum cell sets work: $(11,12,21)$, $(11,12,22)$, $(11,21,22)$, $(12,21,22)$, where $ij=(R_i,C_j)$. Equivalently omit any one cell. Either diagonal orientation gives the same constraint, so orientations are not extra required cases. With two links four strip dots cannot all be joined; three can make a connected tree. Children can compare all two-brace types physically: adjacent braces miss a strip; opposite braces leave two groups. No two works.
\hints{Try removing one brace from D and test.}{Can you find a different empty cell with the same count?}{Compare adjacent and opposite two-brace placements.}
Skip counting every orientation or transcribing a complete catalog. Keep child-selected designs and a contrasting failed attempt. An extension is to choose one brace of a minimum design and predict the effect of removing it: every one is essential.

\newpage
\titleline{Grades 3--5: physical designs and first link contrasts}
\task{6}{5}{Upper-table core}
Minimum \textbf{four} for the 2-by-3 model. One example is $(11,12,13,21)$: all top-row cells and bottom-left. There are 12 minimum cell sets; the task asks for several, not all. An economical full adult characterization: from all six cells omit any pair \emph{except the two cells in one column}. Thus every remaining four-cell set covers the columns and connects both rows through a column having two braces. Two omissions in one column isolate that column. With five strip dots three links cannot suffice.
\hints{Try a design that holds, then remove one brace and test.}{Can you leave one column without a brace?}{Look for a column that joins both row directions.}
Do not supply a placement algorithm first. Children should choose and revise their designs. Further: compare successful designs with the same count rather than require all 12.

\task{7}{6}{Core if children can retain two label systems}
Keep the worked conversion at the top of page 6 before first use. For A draw links $11,12,21,22$; for B draw $11,12,13,21$. \textbf{Circle B only}. Both have four braces. A's group is $\{R1,R2,C1,C2\}$ plus isolated C3, so its right-hand strip can change. B links all five dots. The figures below are answer pictures, not physical pins.
\begin{center}
\begin{tabular}{cc}A: flexible&B: holds\\
\links{2}{3}{1/1,1/2,2/1,2/2}&\links{2}{3}{1/1,1/2,1/3,2/1}
\end{tabular}
\end{center}
\hints{Build both designs and choose where to push.}{Point to one brace's row and column before drawing its link.}{Keep all five dots; which physical direction has no link in A?}
The hand test supports a conjecture. Complete the page-4 two-group handoff before generalizing from pictures alone. If it does not land, physical P6/7 already provide a substantial hour.

\task{8}{7}{Readiness-dependent continuation/return visit}
\textbf{No}: covering every row/column is insufficient. A links $11,12,21,22,33$ and has groups $\{R1,R2,C1,C2\}$ and $\{R3,C3\}$. Its two groups can turn differently, preserving every braced cell. B links $11,12,13,21,31$, connects all six dots and holds. Both have five braces and complete strip coverage. Children draw \emph{both} link pictures before deciding the universal question.
\begin{center}
\begin{tabular}{cc}A: flexible&B: holds\\
\links{3}{3}{1/1,1/2,2/1,2/2,3/3}&\links{3}{3}{1/1,1/2,1/3,2/1,3/1}
\end{tabular}
\end{center}
\hints{Recall the physical covered two-group frame.}{Trace only actual links from a chosen dot; a crossing does not join.}{Which two physical direction groups could take separate small turns?}
If the model explanation remains unclear, save this claim as a conjecture and revisit with the handoff or a rehearsed optional 3-by-3 model. Do not claim the paper design was physically tested. Extension: in A, any empty cell bridges the two groups, so any one new brace makes it hold.

\newpage
\titleline{Grades 3--5: removal and addition answers}
For adult lookup on this page, call the top-row cells $a,b,c$ and the bottom-row cells $d,e,f$, both from left to right. A letter denotes a cell/row--column link, never a physical joint. Student pages use the row/column labels.

\task{9}{8}{Continuation after the handoff}
The starting set is $a,b,c,d,e$ (all except bottom-right $f$). Removable single braces are \textbf{$a,b,d,e$}, cells $(11),(12),(21),(22)$. Mark those four on both pictures. \textbf{$c=(13)$ is essential}: its removal isolates C3. The other four links form a cycle, and each can be removed while a path through the other three remains. This is four \emph{separate} tests, not removal of all four together. Restore the same design before each next trial.
\hints{Choose one brace, remove it, test, then restore it.}{In the link picture, is there another route between that link's ends?}{Which column has only one connection?}
Extension: after a permitted removal, the four remaining links make a tree; every further single removal then destroys rigidity. Ask children to choose the next removal and test that new reference deliberately.

\task{10}{8}{Continuation; can remain wholly physical}
Start afresh with all six cells braced; remove both braces in each chosen pair simultaneously. There are 15 unordered pairs. The \textbf{three failing pairs} are $ad,be,cf$: both cells of one column are removed, isolating that direction. The \textbf{12 successful pairs} are $ab,ac,ae,af,bc,bd,bf,cd,ce,de,df,ef$. Every surviving four-link design is connected. Each starts from the full six-brace reference, not from P9's five-brace frame or the previous trial. Four printed maps are for chosen contrasting trials; children need not reproduce all 15.
\hints{Choose two removals and keep their combined record.}{What happens if both come from the same column?}{If every column remains linked, is there a column that still joins both rows?}
Child explanation: with four braces among three columns, some column has both rows. If no column is empty, each other column links to a row already in that group, so all dots join. Extension: pool independently chosen tests and keep a checkable catalog rather than impose a long individual transcription.

\task{11}{9}{Paper continuation after the handoff}
\textbf{A: no place works.} Its links $11,12,21,22$ form one four-dot group with R3 and C3 separate: three groups. One new link can reduce their number by only one. Adding $33$ joins the two isolated dots but still leaves two groups. Adding a third-row or third-column brace connects only one isolated direction.

\textbf{B: exactly four places work}, $(13),(23),(31),(32)$. Its links $11,12,22$ join $\{R1,R2,C1,C2\}$, and $33$ joins $\{R3,C3\}$. The four listed cells bridge these two groups. The remaining empty cell $(21)$ stays inside the first group and adds a cycle without joining the other group. Mark all four working cells, rather than just one example.
\begin{center}
\begin{tabular}{cc}A: three groups&B: two groups\\
\links{3}{3}{1/1,1/2,2/1,2/2}&\links{3}{3}{1/1,1/2,2/2,3/3}
\end{tabular}
\end{center}
\hints{Trace the separate groups, retaining unused dots.}{Choose an empty cell; point to the two groups its link would join.}{Can one link join three separate groups?}
If children are struggling, use the physical two-group comparison again instead of adult tracing all answers. Extension: A needs two additions; for example $(13)$ and $(31)$ connect all groups. Let children invent other two-addition repairs.

\newpage
\fontsize{10.5}{12.3}\selectfont
\titleline{Ready Grades 4--5: minima and universal claims}
These are ideal complete-grid problems after the physical-to-link handoff, usually a return visit. No 4-by-5 or 4-by-4 physical kit is promised. Children choose a design, counterexample or drawing argument; adults can read/record their explanation.

\task{12}{10}{Ready continuation/return visit}
Minimum \textbf{eight}: there are four row dots and five column dots. Start with nine groups. Each brace joins at most two, decreasing the number by at most one; seven braces cannot leave just one group. An eight-brace example is all five top-row cells plus the first-column cells in rows 2, 3, 4: $11,12,13,14,15,21,31,41$. Those links connect every dot. This proves necessity \emph{and} sufficiency; the count alone is not a test of a particular placement.
\hints{Draw the row and column dots before adding links.}{What is the greatest number of separate groups one new link can join?}{Try making a connected picture with exactly eight links.}
Extension: remove any brace from a minimum design; its tree splits. Adding a ninth creates a cycle and may provide a removal choice. Children can search for this on their own diagram.

\task{13}{11}{Ready universal argument}
\textbf{Yes}, for this precise 3-by-3, six-brace, complete-coverage situation. Suppose its link picture were disconnected. Coverage means every component has at least one row and one column. If it has two groups, their row sizes must be $1+2$ and column sizes $1+2$. Matching large with large allows at most $1\cdot1+2\cdot2=5$ distinct links; matching large with small allows at most $1\cdot2+2\cdot1=4$. If it has three groups, each has one row/column and there are at most three links. Thus six cannot fit without a cross-group link: all dots are connected. One brace per cell matters to this bound.
\hints{Try to build a covered design that stays in two separate groups.}{For each group, count the available row--column pairs inside it.}{Can two such groups contain six different cells?}
A drawing explanation can circle two proposed groups and shade every cell available inside them. Code checked all 84 six-cell subsets: 78 cover all strips, and all 78 are connected. Exhaustive testing is supporting evidence; the capacity argument explains why all must work. Optional next question: what changes with four rows/columns?

\task{14}{12}{Ready counterexample; depth worth preserving}
\textbf{No}. Brace eight cells of the upper-left 3-by-3 block, leaving its $(3,3)$ empty, and also brace $(4,4)$. That is exactly nine: $11,12,13,21,22,23,31,32,44$. Every row/column has a brace. The first six strip dots form one connected group, and $\{R4,C4\}$ is another. All cells crossing between these groups are empty; the groups can take different small turns while preserving the diagonal lengths. This supplies a finite-flex counterexample, not just a brace-count objection.
\begin{center}\grid{4}{4}{13}{1/1,1/2,1/3,2/1,2/2,2/3,3/1,3/2,4/4}\end{center}
\hints{Try dividing rows and columns into a big group and a small group.}{How many brace positions fit inside a 3-row/3-column block?}{Can you use eight there and one in the remaining row/column without linking the groups?}
Extension: with coverage, a disconnected 4-by-4 graph has at most $3\cdot3+1\cdot1=10$ links (other two-group splits yield at most 8, and three/four groups at most 6/4). Eleven braces therefore force connectivity; ten can still fail using all nine cells in that block plus $(4,4)$. Treat this as optional further work, not a requirement to finish the packet.

\newpage
\fontsize{11}{13.6}\selectfont
\titleline{Short proofs for the adults and source record}
\sub{First order and finite motion are separate arguments}
At square joints $p(i,j)=(j,i)$, a velocity $(u,v)$ satisfies horizontal bars only if $u(i,j)=a_i$, and vertical bars only if $v(i,j)=b_j$. Let $A_i=a_i-a_{i-1}$, $B_j=b_j-b_{j-1}$. Either diagonal orientation imposes $A_i+B_j=0$. Thus row turns $-A_i$ equal column turns $B_j$ precisely along links. One turn per component gives $k$ variables, one being whole-frame rotation; two more variables are translations. The rank is $2(m+1)(n+1)-2-k$. This establishes the infinitesimal criterion at the square placement.

For finite motion, write a joint locally as
\[
p(i,j)=p(0,0)+\sum_{t=1}^{j}H_t+\sum_{s=1}^{i}V_s,\qquad |H_t|=|V_s|=L.
\]
Equal opposite cell sides preserve the parallelogram branch near the nondegenerate start. A brace fixes $|H_j\pm V_i|=L\sqrt2$, hence $H_j\cdot V_i=0$. In a connected graph these right angles propagate, so every $H$ is the same rotated horizontal vector and every $V$ the same perpendicular vertical vector. Local motion is only rigid motion, hence no continuous shape-changing path can start here. For disconnected components rotate each component's $H$ and $V$ together by its own small angle. The sums define joint positions with every bar and chosen diagonal length unchanged; a cross-component cell changes angle. This is an actual flex construction. For triangle rigidity, fixing a side restricts the third vertex to at most two circle intersections; it cannot follow a continuous shape-changing path.

\sub{Exact primary mathematical sources}
Bolker and Crapo, \emph{Bracing Rectangular Frameworks. I}, SIAM J. Appl. Math. 36(3), 1979, pp. 473--490, \href{https://epubs.siam.org/doi/abs/10.1137/0136036}{DOI 10.1137/0136036}. The publisher abstract was inspected and identifies minimal planar square-grid bracings with bipartite spanning trees. The full original paper was unavailable; no internal proof/page was claimed read.

Grasegger and Legersk\'y, \emph{Bracing frameworks consisting of parallelograms}, \href{https://arxiv.org/abs/2008.11521v1}{arXiv:2008.11521v1} (2020), Theorem 1.1 p. 3; Definition 3.8 p. 8 and 3.17/3.19 pp. 13--14; finite-flex/graph proofs pp. 16--17. Those pages were read in the local downloaded primary PDF. The paper treats the broader class of braced P-frameworks; this guide uses a separately written elementary proof for complete square grids only. Its p. 17 explicitly distinguishes grids with holes.

\sub{Teaching lessons actually found, and our design inferences}
Rozhkovskaya, \emph{Math Circles for Elementary School Students} (AMS, 2014), Lesson 7, ``At the lesson,'' \S2, printed p. 122 (local EPUB part0017): reports sustained, enthusiastic K'nex polygon construction. Givental, Nemirovskaya and Zakharevich, \emph{Math Circle by the Bay} (AMS, 2018), ``How we teach,'' printed p. ix (PDF p. 10): discusses manipulatives, accessible unambiguous statements, independent solving and practice verbalizing ideas. Neither reports piloting this bracing packet.

Local precedents: year-1 \emph{Handouts 3}, P3.4--3.7, pattern-block hexagon/symmetry work; \emph{Handouts 4}, P4.2--4.5, varied filling and symmetry grouping. The organizer's stronger classroom evidence is the concrete handling versus paper-rule difficulty recorded in the current workflow context. Our preassembled kits, strip dimensions, chosen handoff, routes and timing are \emph{untested design inferences}. Mathematical ideas are established; guide wording, proofs, lookup tables and figures are newly authored. No borrowed exemplar text, prompt bundle, third-party figure or book excerpt belongs in this guide source.

\newpage
\titleline{Return visits, records and verification limits}
\sub{Use the same theme for a deeper investigation}
First visit can end after concrete P1--7, with no graph criterion claimed as understood. Save one brace map per pair/trio and the physical question they wanted to keep trying. On return, reconstruct their own design before new work. K--1 can investigate which different-length diagonal fits a chosen rhombus, with an adult-prepared matching bar; this changes the model and is not a square-grid application. Third graders can invent hard-to-classify 2-by-2/2-by-3 designs and exchange predictions. The upper table can revisit its two-group physical frame, then choose redundancy/repair P9--11 or minimum/universal P12--14. These are different investigations of constraints, not merely larger numbers.

Keep the concrete entry available alongside proofs. For the eleven-child group, adults remain at their tables and grade labels do not change the staffing. If third graders use the spare 2-by-3 kit, make the handoff after the upper table finishes with it, without requiring an adult to supervise another table. A second visit can focus on one substantial proof/counterexample rather than finish every page.

\sub{Write down after the session}
For each work group record date, table, problems actually used and the designs saved. Distinguish ``printed/available'' from ``used.'' Record a child's claim or drawing, what they tried, and whether it was guessed, hand-tested, checked in several cases or explained from fixed lengths. Note whether children understood shape change versus whole-frame motion, brace-to-direction meaning, all dots/crossings, and restore-before-next-test. Record where adults read/serviced pins versus chose placements or traced groups. Log confusing instructions, favorite continuations, waiting for pin service and physical fit problems. Separate observations from hypotheses about why they happened; use the evidence to change pacing or materials next time.

\begin{tabularx}{\linewidth}{@{}lX@{}}\textbf{Group/date} & \rule{.8\linewidth}{.3pt}\\[8pt]
\textbf{Problems used} & \rule{.8\linewidth}{.3pt}\\[8pt]
\textbf{Saved design/claim} & \rule{.8\linewidth}{.3pt}\\[8pt]
\textbf{Evidence status} & guessed / tested / cases checked / explained\\[8pt]
\textbf{Adult help/fit} & \rule{.8\linewidth}{.3pt}\\[8pt]
\textbf{Wanted next} & \rule{.8\linewidth}{.3pt}
\end{tabularx}

\sub{Independent digital checks for this guide}
The authored \texttt{verify\_guide.py} re-computes every stated minimum, printed-design component set, single/pair deletion, addition location, P13 exhaustive result and P14 counterexample without importing student checks. Exact rational rigidity-matrix ranks independently agree with the graph predictions. A small finite-flex vector construction checks unchanged side lengths and both diagonal orientations, and measures a changed unbraced diagonal. K--1 lengths and kit counts are also recalculated. Optional actual-PDF auditing reads brace paths and equal cell axes from all printed braced task grids in the relevant student pages; no source comments are used as answer evidence.

Every guide page is rendered and visually checked. A copied-source clean build compares page count, Letter dimensions, text and rendered pixels from another working directory. These checks establish digital/math consistency, \emph{not} physical fabrication, handling, staffing rehearsal or classroom results. Both physical rehearsal and classroom piloting remain unperformed. The guide is separately authored for F52-S-v2; current digital checks cover this guide. No remote copy has been updated by this work.
\end{document}
